% Part: lambda-calculus
% Chapter: introduction
% Section: church-rosser

\documentclass[../../../include/open-logic-section]{subfiles}

\begin{document}

\olfileid{lam}{int}{cr}
\olsection{丘奇--Rosser性质}

\begin{thm}
\ollabel{thm:church-rosser}
设 $M$、$N_1$ 和 $N_2$ 为项，且满足 $M \red N_1$ 和 $M \red
N_2$。则存在项 $P$，使得 $N_1 \red P$ 且 $N_2 \red P$。
\end{thm}

\begin{cor}
若 $M$ 可归约为范式，则该范式是唯一的。
\end{cor}

\begin{proof}
如果 $M \red N_1$ 且 $M \red N_2$，根据前一定理，存在项~$P$ 使得 $N_1$ 和 $N_2$ 均可归约到~$P$。若 $N_1$ 和~$N_2$ 均为范式，则仅当 $N_1 \ident P \ident
N_2$ 时才可能。
\end{proof}

最后，若两个项 $M$ 和~$N$ 能归约到同一个项，则称它们是\emph{$\beta$-等价的}，或简称为\emph{等价的}；
换言之，即存在某个 $P$ 使得 $M
\red P$ 且 $N \red P$。这记作 $M \equal[\beta] N$。利用
\olref{thm:church-rosser}，你可以验证 $\equal[\beta]$ 是一个等价关系，并且它还满足额外的性质：对任意 $M$
和~$N$，如果 $M \red N$ 或 $N \red M$，则 $M \equal[\beta] N$。（事实上，可以证明 $\equal[\beta]$ 是满足此性质的\emph{最小}等价关系。）

\end{document}